\subsubsection{Método}\label{sec:inicio-metodo}
La fase de Inicio adaptará DAD para reducir las incertidumbres que condicionan la viabilidad del producto antes de ampliar la producción de escenarios. Se combinarán la revisión de necesidades y antecedentes, la especificación de requisitos, el diseño inicial de arquitectura y las prácticas controladas con FDS, de manera que las decisiones se apoyen en entradas y resultados identificables. La adaptación toma como base las metas de proceso que DAD define para la fase de Inicio. Se conservan sin cambios la formación del equipo, el desarrollo de una visión común, la exploración del alcance, la identificación de la estrategia técnica, el plan de entrega y la estrategia de pruebas. Estas metas se concretan en la asignación de roles, la especificación de requisitos según ISO/IEC/IEEE 29148 [9], la arquitectura inicial, una prueba de concepto con FDS y los criterios iniciales de E6. Tres metas se ajustan al carácter académico del proyecto. La alineación con la dirección de la organización se interpreta como alineación con la plantilla de trabajo de grado, el director y la retroalimentación del PAFP [4]. La formación del entorno de trabajo se concreta en el repositorio, la especificación de requisitos, la cuenta en el clúster y las herramientas SIG. Como el proyecto no cuenta con presupuesto ni institución cliente, la meta de asegurar el financiamiento se sustituye por asegurar los recursos de los que depende: acceso al HPC, datos con licencia abierta y una revisión ejecutable de FDS. A estas metas se suma la gestión de riesgos, que DAD trata como meta transversal. Esta adaptación responde a la dependencia del software respecto de datos científicos y recursos de cálculo, por lo que un resultado de fase no será solamente un documento de diseño, sino también evidencia suficiente para preparar y reproducir los casos posteriores. La preparación de esta propuesta constituye un antecedente académico y no sustituye la ejecución de las actividades técnicas que se describen a continuación.
